Tables~\ref{tab:main} and~\ref{tab:rel} give the main results; Fig.~\ref{fig:bars} plots MRR per system and setting. All numbers are means over \rhval{const/n_seeds} seeds; differences quoted below are paired over seeds.

\begin{table*}[t]\centering\caption{Main results: filtered MRR, Hits@1 and Hits@10 (mean $\pm$ std over seeds) on held-out triples of the training graph (transductive) and on a fresh graph with new entities (inductive). Bold: Path-MP; underline: best baseline per column and setting.}\label{tab:main}
\begin{tabular}{llcccc}
\toprule
Method & Task & mrr $\uparrow$ & hits1 $\uparrow$ & hits10 $\uparrow$ & $n$ \\
\midrule
Rule oracle (hand-written) & transductive & \underline{0.937 $\pm$ 0.026} & \underline{0.931 $\pm$ 0.026} & 0.940 $\pm$ 0.026 & 5 \\
TransE+fold-in (reimplemented) & transductive & 0.656 $\pm$ 0.029 & 0.423 $\pm$ 0.056 & 0.981 $\pm$ 0.005 & 5 \\
TransE (reimplemented) & transductive & 0.658 $\pm$ 0.024 & 0.424 $\pm$ 0.049 & \underline{0.983 $\pm$ 0.005} & 5 \\
\textbf{Path-MP (2-hop)} & transductive & \textbf{0.988 $\pm$ 0.004} & \textbf{0.985 $\pm$ 0.003} & \textbf{0.990 $\pm$ 0.005} & 5 \\
\midrule
Rule oracle (hand-written) & inductive & \underline{0.950 $\pm$ 0.020} & \underline{0.946 $\pm$ 0.019} & 0.952 $\pm$ 0.020 & 5 \\
TransE+fold-in (reimplemented) & inductive & 0.644 $\pm$ 0.025 & 0.423 $\pm$ 0.046 & \underline{0.982 $\pm$ 0.010} & 5 \\
TransE (reimplemented) & inductive & 0.011 $\pm$ 0.002 & 0.000 $\pm$ 0.000 & 0.015 $\pm$ 0.006 & 5 \\
\textbf{Path-MP (2-hop)} & inductive & \textbf{0.989 $\pm$ 0.007} & \textbf{0.987 $\pm$ 0.008} & \textbf{0.992 $\pm$ 0.007} & 5 \\
\bottomrule
\end{tabular}

\end{table*}

\begin{table*}[t]\centering\caption{Filtered MRR per relation (mean $\pm$ std over seeds).}\label{tab:rel}
\begin{tabular}{llccccc}
\toprule
Method & Task & mrr\_parent $\uparrow$ & mrr\_sibling $\uparrow$ & mrr\_grandparent $\uparrow$ & mrr\_uncle $\uparrow$ & $n$ \\
\midrule
Rule oracle (hand-written) & transductive & \underline{0.828 $\pm$ 0.035} & \underline{0.976 $\pm$ 0.020} & \underline{0.996 $\pm$ 0.009} & \underline{0.958 $\pm$ 0.036} & 5 \\
TransE+fold-in (reimplemented) & transductive & 0.735 $\pm$ 0.035 & 0.485 $\pm$ 0.010 & 0.858 $\pm$ 0.065 & 0.498 $\pm$ 0.052 & 5 \\
TransE (reimplemented) & transductive & 0.732 $\pm$ 0.045 & 0.485 $\pm$ 0.008 & 0.867 $\pm$ 0.056 & 0.499 $\pm$ 0.041 & 5 \\
\textbf{Path-MP (2-hop)} & transductive & \textbf{0.964 $\pm$ 0.019} & \textbf{0.999 $\pm$ 0.002} & \textbf{0.996 $\pm$ 0.009} & \textbf{0.994 $\pm$ 0.007} & 5 \\
\midrule
Rule oracle (hand-written) & inductive & \underline{0.851 $\pm$ 0.065} & \underline{0.990 $\pm$ 0.002} & \underline{1.000 $\pm$ 0.000} & \underline{0.967 $\pm$ 0.033} & 5 \\
TransE+fold-in (reimplemented) & inductive & 0.692 $\pm$ 0.035 & 0.480 $\pm$ 0.006 & 0.863 $\pm$ 0.048 & 0.506 $\pm$ 0.033 & 5 \\
TransE (reimplemented) & inductive & 0.012 $\pm$ 0.004 & 0.010 $\pm$ 0.005 & 0.011 $\pm$ 0.003 & 0.012 $\pm$ 0.004 & 5 \\
\textbf{Path-MP (2-hop)} & inductive & \textbf{0.960 $\pm$ 0.025} & \textbf{1.000 $\pm$ 0.000} & \textbf{1.000 $\pm$ 0.000} & \textbf{0.998 $\pm$ 0.003} & 5 \\
\bottomrule
\end{tabular}

\end{table*}

\begin{figure}[t]\centering\includegraphics[width=\columnwidth]{bars_main_mrr.pdf}
\caption{Filtered MRR per system, in the two settings (mean and spread over seeds, as drawn by \texttt{rh fig bars}).}\label{fig:bars}\end{figure}

\vspace{4pt}\noindent\textbf{H1: Path-MP beats TransE on held-out triples (supported).} Path-MP reaches MRR \rhval{main/path-mp-2-hop/transductive/mrr/mean:3} against \rhval{main/transe-reimplemented/transductive/mrr/mean:3} for tuned TransE, a paired difference of \rhval{cmp/main/transe-reimplemented/transductive/mrr/delta:3} (paired $p=\rhval{cmp/main/transe-reimplemented/transductive/mrr/paired_p:sci2}$ over \rhval{const/n_seeds} seeds, the test registered in the proposal). The picture is different for Hits@10: TransE reaches \rhval{main/transe-reimplemented/transductive/hits10/mean:3} against \rhval{main/path-mp-2-hop/transductive/hits10/mean:3}, but Hits@1 is \rhval{main/transe-reimplemented/transductive/hits1/mean:3} against \rhval{main/path-mp-2-hop/transductive/hits1/mean:3}. TransE places the answer among the first ten candidates almost always but often not first, which, under a 500-entity candidate set, is a weak criterion; MRR and Hits@1 expose the difference. Path-MP also exceeds the one-step rule oracle (MRR \rhval{main/rule-oracle-hand-written/transductive/mrr/mean:3}, difference \rhval{cmp/main/rule-oracle-hand-written/transductive/mrr/delta:3}, paired $p=\rhval{cmp/main/rule-oracle-hand-written/transductive/mrr/paired_p:sci2}$). The oracle is a one-step reference, not a ceiling: the learned model can combine several rule paths, such as reaching a grandparent through a parent's uncle edge.

\vspace{4pt}\noindent\textbf{Per relation (Table~\ref{tab:rel.}) and a failure case.} TransE is best on grandparent (MRR \rhval{main/transe-reimplemented/transductive/mrr_grandparent/mean:3}) and worst on sibling and uncle (\rhval{main/transe-reimplemented/transductive/mrr_sibling/mean:3} and \rhval{main/transe-reimplemented/transductive/mrr_uncle/mean:3}). A plausible explanation, which we did not isolate with a run, is that sibling is symmetric and a translation model can represent a symmetric relation only with a near-zero relation vector, which makes it hard to separate siblings from other near neighbours; grandparent is a composition of parent, which translations represent naturally. Path-MP is nearly perfect on sibling, grandparent and uncle but weakest on parent (\rhval{main/path-mp-2-hop/transductive/mrr_parent/mean:3} transductive, \rhval{main/path-mp-2-hop/inductive/mrr_parent/mean:3} inductive). The oracle is weakest on parent as well (\rhval{main/rule-oracle-hand-written/transductive/mrr_parent/mean:3} transductive), consistent with a parent triple being recoverable only when a sibling with the same parent is observed, which is not the case for every held-out parent triple; Path-MP's parent MRR is the clearest remaining failure of the model in this world.

\vspace{4pt}\noindent\textbf{H2: TransE collapses on new entities, Path-MP does not (supported).} Frozen TransE gets MRR \rhval{main/transe-reimplemented/inductive/mrr/mean:3} on the new graph and Hits@1 \rhval{main/transe-reimplemented/inductive/hits1/mean:3}, a change of \rhval{analysis/analysis-final/inductive/shift_te_diff/mean:3} relative to the transductive setting (paired $p=\rhval{analysis/analysis-final/inductive/shift_te_p/mean:sci2}$). This is the expected result, since the vectors of the new entities are random; it is included as the reference point for the fold-in baseline. Path-MP's change from transductive to inductive is \rhval{analysis/analysis-final/inductive/shift_pm_diff/mean:3} MRR (paired $p=\rhval{analysis/analysis-final/inductive/shift_pm_p/mean:2}$): within the 0.05 tolerance stated in H2, and statistically indistinguishable from zero with five seeds. The new graph is generated by the same process as the training graph, so this tests transfer to new entities, not to a different distribution of graphs.

\vspace{4pt}\noindent\textbf{H3: fold-in recovers most of TransE's loss but stays below Path-MP (supported).} TransE with fold-in reaches \rhval{main/transe+fold-in-reimplemented/inductive/mrr/mean:3} on the new graph, \rhval{analysis/analysis-final/inductive/h3_tf_vs_te_ind_diff/mean:3} above frozen TransE (paired $p=\rhval{analysis/analysis-final/inductive/h3_tf_vs_te_ind_p/mean:sci2}$), and its inductive MRR is not distinguishable from TransE's transductive MRR (change \rhval{analysis/analysis-final/inductive/shift_tf_diff/mean:3}, $p=\rhval{analysis/analysis-final/inductive/shift_tf_p/mean:2}$). Path-MP stays above fold-in by \rhval{analysis/analysis-final/inductive/h3_pm_vs_tf_ind_diff/mean:3} MRR (paired $p=\rhval{analysis/analysis-final/inductive/h3_pm_vs_tf_ind_p/mean:sci2}$). The conclusion that TransE cannot handle new entities is thus an artefact of the evaluation set-up: with fold-in, the inductive handicap of TransE vanishes in this world, and what remains is its gap to the path model in both settings. Fold-in costs training on the new graph at test time, while Path-MP only needs a forward pass.
